\section{Introduction}
\label{sec:intro}

Enterprises record regulations, failure analyses and operational lessons in
unstructured documents. A document collection becomes a knowledge base only
when its documents are parsed and admitted reliably, classified into the
collections that own them and kept consistent as those collections change.
A question must then be answerable by moving down a path from the collection
to the document and finally to the passage that carries the evidence.
Answering a question such as \emph{why did edge delamination appear in a
photovoltaic inspection?} requires the collection that holds inspection
reports, the report that records the decision, the passage inside it and an
accurate pointer to both.

The two capabilities required for this operate on different parts of the
lifecycle. Modular RAG and agentic-retrieval surveys already treat indexing
and pre-retrieval as first-class pipeline modules
\cite{modular-rag,gao-ragsurvey}; what remains separated is the governance of
the collection itself---its placement, lifecycle and per-tier audit---from
the retrieval step that consumes it. Content-based
retrieval addresses the query step within a collection that is already
indexed: retrieval-augmented generation (RAG) embeds queries and passages in
a shared vector space \cite{rag-lewis,dpr}, agentic systems let a language
model call retrieval tools \cite{mcp,singh-agentic-rag}, and hybrid, lexical
and graph-structured variants extend the ranking model
\cite{bm25,rrf-cormack,graphrag,lightrag,raptor}. Knowledge management
addresses the collection itself: enterprise content management governs
organisational documents \cite{paivarinta-ecm-2005,tyrvaainen-ecm-2006},
metadata and tagging studies ask how documents are described and reviewed
\cite{park-metadata-quality-2009,golder-tagging-2006}, and public knowledge
bases curate facts under explicit ontologies
\cite{auer-dbpedia-kb-2007,hoffart-yago2-kb-2013}. Semantic similarity
measures topical relatedness rather than evidence containment, a failure mode
discussed in the hallucination and RAG surveys \cite{siren,gao-ragsurvey}.
A correctly ranked passage is
likewise unreachable when the collection that owns it was never classified, was filed
under a name that does not match its content, or was indexed after a rename
that left the index stale.

This separation creates a systems problem at the boundary between the two
streams. A collection can be perfectly indexed and the system still answer wrongly,
because the first decision a question requires---which collection to
consult---is made without reference to the collection structure in flat,
single-index deployments.
Distributed IR has long studied this collection-selection step
\cite{callan1995searching,si2003redde}, but models collections as static
statistical summaries rather than as managed state that is written to and
checked by reconciliation audits after selected lifecycle operations.
Conversely, a carefully curated hierarchy does not help when no tiered
procedure carries a question from the catalogue down to the evidence. The
question we study is therefore not how to rank passages better, but how a
document collection should be \emph{managed} so that a query can be served
tier by tier---selecting a collection, then a document, then a segment---and
so that each tier can be audited independently. We study that question through
an implemented system and bounded operational audits. Integration by itself is
not treated as evidence of improved retrieval or answer quality.

HMRS (Hierarchical Management Retrieval System) is a self-hosted
knowledge document management system. It maintains collections, documents and
segments and supports a catalogue--shelf--document--segment access path.
An ingest agent selects collections from parsed content, while lifecycle audits
compare relevant files and registries. The optional evidence gate scores
candidate segments; it does not prove support. Answer workflows distinguish a
completed search without evidence from tool failure. These stages are described
in Sections~\ref{sec:system}--\ref{sec:implementation}.

The platform integrates document parsing, automatic collection assignment,
lifecycle management, operational experiences, the tiered access path and a
Model Context Protocol (MCP) server \cite{mcp}. It
runs locally and exposes the management and retrieval functions to several
agent harnesses. We distinguish implemented functions from functions exercised
by the reported campaigns.

\noindent\textbf{Contributions.}
\begin{enumerate}
\item A management-level model for knowledge document management and its
implementation: documents are managed at collection, document and segment
level, collections are assigned automatically from parsed content,
sub-libraries are nodes with their own indices, and the state of a
collection is described by three sources that can be compared in lifecycle
audits. The system description makes explicit which properties
are enforced by tools and which remain agent-workflow decisions
(Sections~\ref{sec:system}--\ref{sec:implementation}).
\item Collection-level evidence. A 100-paper campaign produced 322 stored parts
in five libraries, of which 79 were retrievable and 21 were not with
incomplete failure attribution; a separate cross-lingual test measured
automatic membership accuracy of 0.875 for English against 1.0 for Chinese and
Japanese; and a 72-question ten-domain collection-selection audit returned the
gold knowledge base for 72/72 questions in one arm and nominally name-only for
37/72 in another, where an unlogged fallback could still use descriptions, so
the contrast is a configuration comparison and not an isolated treatment effect
(Section~\ref{sec:routing-audit}).
\item Behaviour of the lower tiers of the access path. A mechanical full-read
script found the target strings in 72/72 cases after reading 879 parts
(24.3 MB, 253.7 s per question) but generated no answers, and two ten-question
answer grids reported gold identifiers in the answer or recorded trace for all 60 cells while strict and
alias-aware lexical coverage and recognised citation markers vary with the
corpus substrate. These deterministic checks describe answer-layer behaviour
and do not establish semantic correctness or claim-to-source support
(Section~\ref{sec:e5}, Table~\ref{tab:e5-answers}). A same-corpus content-layer
contrast further shows the tiered path with a chunk-window evidence union
containing the gold string in 52/72 questions, above a scoped flat union
(39/72) and a precise per-document arm without navigation (46/72). The arms
differ in several components, and a gold-label-dependent lookup limits the
flat comparison
(Section~\ref{sec:ladder-opt}).
\item A same-pool diagnostic on the segment tier: score-based ordering lowers
Hit@5 from 0.94 to 0.44, while fusion order among recorded survivors restores
0.94, and all 191 recorded scores exceed the configured 0.5 threshold, so the
measurable risk lies in ordering; filtering was not activated on this pool
and therefore remains untested
(Sections~\ref{sec:lanes} and~\ref{sec:e4}--\ref{sec:e6}).
\end{enumerate}

\subsection{Research questions and reading guide}
\label{sec:rq}

Five questions connect the management model to the evidence and keep each
conclusion within the scope of its measurement. \textbf{RQ1} asks whether
documents can be parsed, automatically classified into a collection and stored
with state that stays consistent across the three sources, and where that
process fails. \textbf{RQ2} asks whether the tiered path selects the correct
collection and narrows to the documents that contain the target evidence, and
what the collection-selection contrast does and does not establish.
\textbf{RQ3} compares the tested retrieval configurations on document recall,
ranking and retrieval cost under the evaluated corpus and substrate
conditions, and asks whether the tiered path contains more gold evidence than
flat or navigation-free alternatives on the same corpus. \textbf{RQ4} examines
how verifier-score ordering affects recall and what can be inferred about
filtering when the configured threshold rejects no recorded score.
\textbf{RQ5} summarises observed abstention, runtime, tool cost and execution
failures at the answer layer.

Section~\ref{sec:related} positions the work; Sections~\ref{sec:system}--
\ref{sec:organization} describe the management model, the collection and
document levels, and the tiered access path;
Section~\ref{sec:implementation} covers deployment and enforcement; and
Section~\ref{sec:eval} evaluates the system, followed by implications,
limitations and conclusions in
Sections~\ref{sec:discussion}--\ref{sec:conclusion}.
